\clearpage
\ExposeTOC
\begin{center}
{\large\textbf{SGA 7}}\\[1.5em]
{\Large\underline{第 XI 讲}}\\[2.5em]
{\Large\underline{完全交的上同调}}\\[1.5em]
{\large\underline{P. DELIGNE}}
\end{center}

\bigskip
\noindent\underline{1. 射影空间与完全交的上同调}

设 $X$ 为纯维数为 $n$ 的非奇异复射影簇, $\underline O(1)$ 为 $X$ 上的丰沛可逆层. 根据 Akizuki 与 Nakano [1], Kodaira 消失定理 [4] 蕴含

\begin{equation}\tag{*}
H^i(X,\Omega_X^j(-1))=0\qquad\text{当 }i+j<n.
\end{equation}

仍由 [1], 此结果蕴含复上同调中的``弱 Lefschetz 定理''. 我们拟用直接计算证明: 若 $X$ 是任意特征的域上的射影空间(projective space), 则 (*) 仍成立. 这样, 遵循 [1], 就能把任意特征的域上的完全交(complete intersection)的 Hodge 上同调和 de Rham 上同调与射影空间的相应上同调加以比较.

然后我们将详细说明所得结果的一些推论. 在 §2 中, 我们给出由 Hirzebruch 的数值公式 (2.2) 推出的定性结果. 感谢 N. Katz 在准备该节时给予的帮助.

1.0. 设 $\underline E$ 为概形 $S$ 上秩为 $r+1$ 的局部自由模. \underline{假设 $r>0$.} 回顾 $\mathbb P(\underline E)$ 表示 $\underline E$ 的超平面所组成的 $S$ 上相对射影空间. 设 $p:\mathbb P(\underline E)\to S$ 为投影, 并令

\begin{equation*}
\eta\in H^0\bigl(S,R^1p_*\Omega^1_{\mathbb P(\underline E)/S}\bigr)
\end{equation*}

为可逆层 $\underline O(1)$ 在 Hodge 意义下的第一 Chern 类.

%%%%%%%%%% source scan idx 47 | folio 40 %%%%%%%%%%

回顾 $\eta$ 是 $\underline O(1)$ 在 $H^0(S,R^1p_*\underline O^*)$ 中的类经对数微分

\begin{equation*}
df/f:\underline O^*\longrightarrow
\Omega^1_{\mathbb P(\underline E)/S}.
\end{equation*}

对于 $\mathbb P(\underline E)$ 上任意 $\underline O$-模层 $\underline F$, 令

\begin{equation*}
\underline F(n)=\underline F\otimes\underline O(1)^{\otimes n}
\qquad(n\in\mathbb Z).
\end{equation*}

\underline{定理} 1.1. (i) \underline{层 $R^ip_*(\Omega^j_{\mathbb P(\underline E)/S}(n))$ 局部自由, 且其构造与任意基变换相容.}

(ii) \underline{对于 $0\leq i\leq r$, $R^ip_*(\Omega^i_{\mathbb P(\underline E)/S})$ 的截面 $\eta^i$ 定义了从 $\underline O_S$ 到此层的同构. 此外,}

\begin{equation*}
R^ip_*\Omega^j_{\mathbb P(\underline E)/S}=0
\qquad\underline{\text{若 }}i\ne j\ \underline{\text{ 或 }}\ i>r.
\end{equation*}

(iii) \underline{若 $n\ne0$, 则除下列两种情形外, $R^ip_*(\Omega^j_{\mathbb P(\underline E)/S}(n))$ 均为零:}

\begin{align}
i&=0\quad\underline{\text{且}}\quad n>j,
  \tag{iii$_a$}\\
i&=r\quad\underline{\text{且}}\quad n<j-r.
  \tag{iii$_b$}
\end{align}

设

\begin{equation*}
\mathbb V(\underline E)=\operatorname{Spec}
\bigl(\operatorname{Sym}^*\underline E\bigr)
\end{equation*}

为 $\underline E$ 定义的仿射空间, $e$ 为其零截面, 并令 $\mathbb V(\underline E)^*=\mathbb V(\underline E)\setminus e(S)$. 设 $a$ 和 $a'$ 分别为 $\mathbb V(\underline E)$ 和 $\mathbb V(\underline E)^*$ 到 $S$ 的投影, $q$ 为 $\mathbb V(\underline E)^*$ 到 $\mathbb P(\underline E)$ 的投影:

\begin{center}
\begin{tikzcd}[row sep=large,column sep=large]
\mathbb V(\underline E)
  \arrow[dr,"a"]
& \mathbb V(\underline E)^*
  \arrow[l,hook']
  \arrow[r,"q"]
  \arrow[d,"a'"']
& \mathbb P(\underline E)
  \arrow[dl,"p"]
\\
{} & S\arrow[ul,bend right,"e"'] & {}
\end{tikzcd}
\qquad .
\end{center}

$\mathbb P(\underline E)$ 上的去零截面仿射空间 $\mathbb V(\underline O(1))^*$ 与 $\mathbb V(\underline E)^*$ 典范同构. 因此, 对于 $\mathbb P(\underline E)$ 上任意拟凝聚 $\underline O$-模层 $\underline F$, 存在典范同构

%%%%%%%%%% source scan idx 48 | folio 41 %%%%%%%%%%

\begin{equation}\tag{1.1.1}
q_*q^*\underline F\simeq
\bigoplus_{n\in\mathbb Z}\underline F(n).
\end{equation}

$q_*q^*\underline F$ 的这个分次也可以由 $\mathbb V(\underline E)$ 的位似群 $\mathbb G_m$ 在此层上的作用得到.

$\mathbb V(\underline E)$ 上的相对向量场可等同于 $\mathbb V(\underline E)$ 到自身的 $S$-态射, 其中将目标视为原点处的切空间. 由 $\mathbb V(\underline E)$ 的恒等映射定义的向量场 $X$ 在位似下不变. 它与 $q$ 的纤维相切, 并定义了一个平凡化

\begin{equation}\tag{1.1.2}
X:\Omega^1_{\mathbb V(\underline E)^*/\mathbb P(\underline E)}
\longrightarrow\underline O_{\mathbb V(\underline E)^*}.
\end{equation}

若 $x_i$ 构成 $\underline E$ 的一组基, 则

\begin{equation*}
X=\sum_i x_i\,dx^i.
\end{equation*}

由于 $\Omega^1_{\mathbb V(\underline E)^*/\mathbb P(\underline E)}$ 的秩为一, 正合列

\begin{equation*}
0\longrightarrow q^*\Omega^1_{\mathbb P(\underline E)/S}
\longrightarrow\Omega^1_{\mathbb V(\underline E)^*/S}
\longrightarrow
\Omega^1_{\mathbb V(\underline E)^*/\mathbb P(\underline E)}
\longrightarrow0
\end{equation*}

诱导正合列

\begin{equation*}
0\longrightarrow q^*\Omega^i_{\mathbb P(\underline E)/S}
\longrightarrow\Omega^i_{\mathbb V(\underline E)^*/S}
\longrightarrow
\Omega^1_{\mathbb V(\underline E)^*/\mathbb P(\underline E)}
\otimes q^*\Omega^{i-1}_{\mathbb P(\underline E)/S}
\longrightarrow0.
\end{equation*}

利用 (1.1.2), 它们变为

\begin{equation}\tag{1.1.3}
0\longrightarrow q^*\Omega^i_{\mathbb P(\underline E)/S}
\longrightarrow\Omega^i_{\mathbb V(\underline E)^*/S}
\xrightarrow{\ \iota_X\ }
q^*\Omega^{i-1}_{\mathbb P(\underline E)/S}
\longrightarrow0.
\end{equation}

典范地有

\begin{equation}\tag{1.1.4}
\Omega^i_{\mathbb V(\underline E)/S}
=a^*\bigwedge^i\underline E.
\end{equation}

所以短正合列 (1.1.3) 由 $\mathbb V(\underline E)^*$ 上的下列长正合列分解而成:

\begin{equation}\tag{1.1.5}
\cdots\longrightarrow {a'}^*\bigwedge^i\underline E
\xrightarrow{\ \iota_X\ }
{a'}^*\bigwedge^{i-1}\underline E
\xrightarrow{\ \iota_X\ }
{a'}^*\bigwedge^{i-2}\underline E\longrightarrow\cdots.
\end{equation}

